\section*{Overview}

Meet-in-the-middle is the standard response to ``the search space is exponential, but \(n\) is only around 40.'' It
does not remove exponential behavior. It cuts the exponent in half and then combines the two halves intelligently.

\section*{Problem-Driven Motivation}

Suppose a subset problem has \(n = 40\). A full brute force explores \(2^{40}\) states, which is completely out of
reach. But \(2^{20}\) is only about one million. That difference is the entire point of the technique:

\[
2^{40} \text{ is impossible},\qquad
2^{20} + 2^{20} \text{ is routine}.
\]

This shows up in subset sum, balanced partition, small-set knapsack, and other ``combine one choice from the left half
with one choice from the right half'' problems.

\section*{Recognition Pattern}

Meet-in-the-middle is a strong candidate when:

\begin{itemize}
  \item \(n\) is roughly \(30\) to \(46\),
  \item the direct search is exponential in \(n\),
  \item a full solution decomposes into an independent left-half choice and right-half choice,
  \item DP by total sum or total value is impossible because the numeric limit is too large.
\end{itemize}

The phrase I look for is:

\begin{quote}
Can I summarize each half independently, then search for a compatible partner?
\end{quote}

\section*{Derivation}

Split the items into two halves:
\[
L = \{a_0,\dots,a_{m-1}\},\qquad
R = \{a_m,\dots,a_{n-1}\}.
\]

Enumerate every subset summary of \(L\) and every subset summary of \(R\). The summary could be:

\begin{itemize}
  \item subset sum,
  \item weight-value pair,
  \item mask property,
  \item any compressed state that combines cleanly.
\end{itemize}

Once those lists are built, the problem is no longer ``search over all \(2^n\) subsets.'' It is ``combine one left
summary with one right summary.''

The merge step is where the real structure of the problem appears:

\begin{itemize}
  \item sort + binary search for subset sum bounds,
  \item two pointers for monotone pairing,
  \item hash map for exact complements,
  \item dominance pruning for multi-attribute summaries.
\end{itemize}

\section*{Worked Problem}

\paragraph{Problem.}
Given \(n \le 40\) positive integers and a limit \(S\), find the largest subset sum not exceeding \(S\).

\paragraph{Why the naive solution fails.}
Enumerating all subsets is \(O(2^n)\), which is too large for \(n=40\).

\paragraph{Step 1: split and enumerate.}
Generate all subset sums of the left half and all subset sums of the right half.

\paragraph{Step 2: sort one side.}
Sort all right-half sums.

\paragraph{Step 3: combine.}
For every left-half sum \(x\), the best partner is the largest right-half sum \(y\) satisfying
\[
x + y \le S.
\]
That partner is found by \texttt{upper\_bound}.

\paragraph{Why this covers all subsets.}
Every full subset is uniquely the union of:

\begin{itemize}
  \item one subset of the left half,
  \item one subset of the right half.
\end{itemize}

So checking all left summaries against all compatible right summaries explores the entire solution space.

\section*{Implementation Reasoning}

The main implementation decisions are:

\begin{itemize}
  \item store all half sums explicitly,
  \item reserve memory for \(2^{n/2}\) elements to avoid repeated reallocations,
  \item use sorting and binary search because the merge condition is an upper bound.
\end{itemize}

In more complex problems, the important refinement is usually \emph{pruning}. If one summary dominates another in all
relevant coordinates, the dominated summary can be discarded before the merge.

\section*{Correctness Intuition}

The split does not approximate anything. It is exact. The only reason it is faster is that combining two half-spaces is
much cheaper than enumerating the full space directly.

\section*{Complexity Analysis}

For the subset-sum version:

\begin{itemize}
  \item enumerate both halves: \(O(2^{n/2})\),
  \item sort one half: \(O(2^{n/2}\log 2^{n/2})\),
  \item binary-search combine: \(O(2^{n/2}\log 2^{n/2})\),
  \item memory: \(O(2^{n/2})\).
\end{itemize}

\section*{Common Pitfalls}

\begin{itemize}
  \item Using meet-in-the-middle when a simpler pseudo-polynomial DP is already better.
  \item Forgetting the memory cost of storing every half summary.
  \item Generating the half summaries correctly but combining them with the wrong inequality or off-by-one boundary.
  \item Missing a dominance-pruning step in a two-dimensional summary problem.
\end{itemize}

\section*{Variants and Failure Modes}

\begin{itemize}
  \item Exact subset-sum queries usually replace sorting with hashing by complement.
  \item Two-attribute problems often sort one side and keep only Pareto-optimal summaries.
  \item If \(n\) is much larger than 45, halving the exponent is still not enough.
\end{itemize}

\section*{Practice Problems}

\begin{itemize}
  \item Maximum subset sum not exceeding a limit.
  \item Closest subset sum to a target.
  \item Pairing one left-half state with one right-half state under a compatibility constraint.
\end{itemize}

\section*{References}

\begin{itemize}
  \item \href{https://usaco.guide/CPH.pdf}{Competitive Programmer's Handbook}.
  \item \href{https://oi-wiki.org/search/mitm/}{OI Wiki: Meet in the Middle}.
  \item \href{https://codeforces.com/catalog}{Codeforces Catalog}.
\end{itemize}
